\documentclass[11pt]{article}

% ============================================================================
% PACKAGES
% ============================================================================
\usepackage[margin=1in]{geometry}
\usepackage{amsmath,amssymb,amsthm}
\usepackage{graphicx}
\usepackage{hyperref}
\usepackage{booktabs}
\usepackage{siunitx}
\usepackage[version=4]{mhchem}
\usepackage{fancyhdr}
\usepackage{enumitem}
\usepackage{mathtools}
\usepackage{braket}
\usepackage{xcolor}
\usepackage{framed}
\usepackage{tcolorbox}
\usepackage{physics}

% ============================================================================
% PAGE SETUP
% ============================================================================
\pagestyle{fancy}
\fancyhf{}
\fancyhead[L]{\small CHM328 Companion Notes -- Lecture 1}
\fancyhead[R]{\small Alastair J.\ Price}
\fancyfoot[C]{\thepage}
\renewcommand{\headrulewidth}{0.4pt}

\hypersetup{
    colorlinks=true,
    linkcolor=blue!60!black,
    citecolor=blue!60!black,
    urlcolor=blue!60!black,
}

% Number equations within sections
\numberwithin{equation}{section}

\title{\textbf{Density Functional Theory:\\From the Hamiltonian to Thermodynamic Properties}\\[1em]
\large Companion Notes for CHM328H1\,S -- Modern Physical Chemistry\\
University of Toronto, Winter 2026}
\author{Alastair J.\ Price\\
\textit{Guest Lecturer}\\
\small Department of Chemistry, University of Toronto}
\date{March 2026}

\begin{document}
\maketitle
\thispagestyle{fancy}

\begin{abstract}
\noindent
These notes accompany the first of two guest lectures on Density Functional Theory delivered for Professor von Lilienfeld's CHM328H1\,S course. The lecture covers the molecular Hamiltonian, the Born--Oppenheimer approximation, Hartree--Fock theory, Kohn--Sham density functional theory, basis sets, the potential energy surface, the harmonic approximation, and the statistical mechanical machinery that connects electronic structure calculations to thermodynamic observables. The central result is the expression $G = E_\text{elec} + E_\text{ZPE} + H_\text{thermal} - TS$, which enables the computation of reaction free energies, equilibrium constants, and related quantities entirely from first principles.
\end{abstract}

\tableofcontents
\newpage

% ============================================================================
\section{Introduction}
\label{sec:intro}
% ============================================================================

Computational chemistry occupies a unique position in modern chemical research: it is both a tool for interpreting experimental observations and a predictive framework in its own right. Where experiment provides direct access to macroscopic observables---rate constants, equilibrium constants, spectroscopic frequencies---computation provides access to the underlying microscopic quantities that determine them. Transition state geometries, reaction barrier heights, individual energy contributions to binding, conformational preferences in flexible molecules---these are quantities that theory can provide and that experiment often cannot measure directly.

The connection between microscopic quantum mechanics and macroscopic thermodynamics is precisely the subject of this course. Statistical mechanics provides the bridge, and the partition function is its central object. What electronic structure theory---in particular, density functional theory (DFT)---supplies is the set of molecular energy levels that enter the partition function. A single DFT frequency calculation provides the electronic energy, the molecular geometry, and the vibrational frequencies; from these, the translational, rotational, vibrational, and electronic partition functions follow, and with them all standard thermodynamic quantities: $\Delta G$, $\Delta H$, $\Delta S$, $K_\text{eq}$, and---as developed in the second lecture---rate constants.

The purpose of these notes is to develop this connection rigorously. We begin with the full molecular Hamiltonian, apply the Born--Oppenheimer approximation to separate the electronic and nuclear problems, survey the principal electronic structure methods (Hartree--Fock and DFT), and then turn to the harmonic approximation and the statistical mechanical framework that converts computed energy levels into thermodynamic properties.


% ============================================================================
\section{The Molecular Hamiltonian}
\label{sec:hamiltonian}
% ============================================================================

\subsection{The Full Hamiltonian}

The starting point of all molecular quantum mechanics is the time-independent Schr\"{o}dinger equation,
\begin{equation}
\hat{H}\Psi = E\Psi,
\label{eq:schrodinger}
\end{equation}
where $\hat{H}$ is the Hamiltonian operator, $\Psi$ is the wavefunction, and $E$ is the energy eigenvalue. For a molecule with $M$ nuclei (with charges $Z_A$ and masses $M_A$) and $N$ electrons, the non-relativistic Hamiltonian in SI units contains five terms:
\begin{equation}
\hat{H} = \underbrace{-\sum_{A=1}^{M}\frac{\hbar^2}{2M_A}\nabla_A^2}_{\hat{T}_\text{nuc}}
\underbrace{- \sum_{i=1}^{N}\frac{\hbar^2}{2m_e}\nabla_i^2}_{\hat{T}_\text{elec}}
+ \underbrace{\sum_{A>B}^{M}\frac{Z_A Z_B e^2}{4\pi\epsilon_0\, r_{AB}}}_{V_\text{nuc-nuc}}
\underbrace{- \sum_{i=1}^{N}\sum_{A=1}^{M}\frac{Z_A e^2}{4\pi\epsilon_0\, r_{iA}}}_{V_\text{elec-nuc}}
+ \underbrace{\sum_{i>j}^{N}\frac{e^2}{4\pi\epsilon_0\, r_{ij}}}_{V_\text{elec-elec}}.
\label{eq:full_hamiltonian}
\end{equation}

The five terms are:
\begin{enumerate}
\item $\hat{T}_\text{nuc}$: nuclear kinetic energy,
\item $\hat{T}_\text{elec}$: electron kinetic energy,
\item $V_\text{nuc-nuc}$: nuclear--nuclear Coulomb repulsion,
\item $V_\text{elec-nuc}$: electron--nuclear Coulomb attraction,
\item $V_\text{elec-elec}$: electron--electron Coulomb repulsion.
\end{enumerate}

The final term, $V_\text{elec-elec}$, is the source of nearly all difficulty in quantum chemistry. It couples all electrons together and prevents the Schr\"{o}dinger equation from being separated into independent one-electron problems. For any system with more than one electron, an exact analytical solution does not exist.

\subsection{Atomic Units}

In atomic units (a.u.), the fundamental constants $\hbar$, $m_e$, $e$, and $4\pi\epsilon_0$ are all set equal to unity. The Hamiltonian simplifies to
\begin{equation}
\hat{H} = -\sum_{A=1}^{M}\frac{1}{2M_A}\nabla_A^2 - \frac{1}{2}\sum_{i=1}^{N}\nabla_i^2 + \sum_{A>B}^{M}\frac{Z_A Z_B}{R_{AB}} - \sum_{i=1}^{N}\sum_{A=1}^{M}\frac{Z_A}{r_{iA}} + \sum_{i>j}^{N}\frac{1}{r_{ij}}.
\label{eq:hamiltonian_au}
\end{equation}
The unit of length is the bohr ($a_0 = \SI{0.5292}{\angstrom}$), and the unit of energy is the hartree ($E_h = \SI{627.5}{kcal/mol} = \SI{27.211}{eV} = \SI{4.360e-18}{J}$). The conversion $\SI{1}{E_h} = \SI{627.5}{kcal/mol}$ is worth memorizing; it appears constantly in computational chemistry. All equations in the remainder of these notes use atomic units unless otherwise stated.

\subsection{The Born--Oppenheimer Approximation}
\label{sec:BO}

The Born--Oppenheimer (BO) approximation~\cite{BornOppenheimer1927} is arguably the most important approximation in all of chemistry. It exploits the mass disparity between nuclei and electrons: a proton is 1836 times heavier than an electron, and heavier nuclei have even larger mass ratios. Consequently, electrons respond effectively instantaneously to nuclear motion. The total wavefunction is approximated as a product,
\begin{equation}
\Psi_\text{total}(\mathbf{r}, \mathbf{R}) \approx \psi_\text{elec}(\mathbf{r}; \mathbf{R})\,\chi_\text{nuc}(\mathbf{R}),
\label{eq:BO}
\end{equation}
where $\mathbf{r}$ denotes all electronic coordinates and $\mathbf{R}$ all nuclear coordinates. The electronic wavefunction $\psi_\text{elec}$ depends parametrically on $\mathbf{R}$---the nuclear positions appear as parameters, not as dynamical variables.

The computational procedure is:
\begin{enumerate}
\item Fix the nuclear positions $\mathbf{R}$.
\item Solve the electronic Schr\"{o}dinger equation at those positions:
\begin{equation}
\hat{H}_\text{elec}\,\psi_\text{elec}(\mathbf{r};\mathbf{R}) = E_\text{elec}(\mathbf{R})\,\psi_\text{elec}(\mathbf{r};\mathbf{R}).
\label{eq:elec_schrodinger}
\end{equation}
\item The electronic energy $E_\text{elec}(\mathbf{R})$ depends on the nuclear positions; this dependence defines the \emph{potential energy surface} (PES).
\end{enumerate}

The electronic Hamiltonian contains three of the five terms from \eqref{eq:hamiltonian_au}:
\begin{equation}
\hat{H}_\text{elec} = -\frac{1}{2}\sum_{i=1}^{N}\nabla_i^2 - \sum_{i=1}^{N}\sum_{A=1}^{M}\frac{Z_A}{r_{iA}} + \sum_{i>j}^{N}\frac{1}{r_{ij}}.
\label{eq:H_elec}
\end{equation}
The nuclear--nuclear repulsion $V_{nn} = \sum_{A>B}Z_AZ_B/R_{AB}$ is a constant for fixed $\mathbf{R}$ and is added at the end:
\begin{equation}
E_\text{total}(\mathbf{R}) = E_\text{elec}(\mathbf{R}) + V_{nn}(\mathbf{R}).
\end{equation}

This---given $\mathbf{R}$ and $N$, find $E_\text{elec}$ and $\psi_\text{elec}$---is the electronic structure problem. It is what Hartree--Fock theory, DFT, and all post-Hartree--Fock methods solve.


% ============================================================================
\section{Hartree--Fock Theory}
\label{sec:HF}
% ============================================================================

\subsection{The Slater Determinant}

The simplest physically meaningful approximation to the many-electron wavefunction is the \emph{Slater determinant}, an antisymmetrized product of $N$ one-electron functions (spin-orbitals) $\{\phi_i\}$:
\begin{equation}
\Psi_\text{HF} = \frac{1}{\sqrt{N!}}
\begin{vmatrix}
\phi_1(\mathbf{x}_1) & \phi_2(\mathbf{x}_1) & \cdots & \phi_N(\mathbf{x}_1) \\
\phi_1(\mathbf{x}_2) & \phi_2(\mathbf{x}_2) & \cdots & \phi_N(\mathbf{x}_2) \\
\vdots & \vdots & \ddots & \vdots \\
\phi_1(\mathbf{x}_N) & \phi_2(\mathbf{x}_N) & \cdots & \phi_N(\mathbf{x}_N)
\end{vmatrix},
\label{eq:slater}
\end{equation}
where $\mathbf{x}_i = (\mathbf{r}_i, \sigma_i)$ includes both spatial and spin coordinates. The determinantal form automatically satisfies the Pauli exclusion principle: interchanging two electrons (rows) changes the sign of $\Psi$, and if two electrons occupy the same spin-orbital, two rows become identical, giving $\Psi = 0$.

\subsection{The Hartree--Fock Energy}

Minimizing $\braket{\Psi_\text{HF}|\hat{H}|\Psi_\text{HF}}$ with respect to the orbitals, subject to orthonormality constraints, yields the Hartree--Fock energy:
\begin{equation}
E_\text{HF} = \sum_{i=1}^{N} h_i + \frac{1}{2}\sum_{i,j=1}^{N}\left(J_{ij} - K_{ij}\right) + V_{nn},
\label{eq:EHF}
\end{equation}
where:
\begin{itemize}
\item $h_i = \braket{\phi_i|-\tfrac{1}{2}\nabla^2 - \sum_A Z_A/r_{iA}|\phi_i}$ are one-electron integrals capturing kinetic energy and nuclear attraction,
\item $J_{ij} = \iint \frac{|\phi_i(\mathbf{r}_1)|^2\,|\phi_j(\mathbf{r}_2)|^2}{r_{12}}\,d\mathbf{r}_1\,d\mathbf{r}_2$ are the \emph{Coulomb integrals}, representing classical electrostatic repulsion between electron clouds,
\item $K_{ij} = \iint \frac{\phi_i^*(\mathbf{r}_1)\phi_j(\mathbf{r}_1)\,\phi_j^*(\mathbf{r}_2)\phi_i(\mathbf{r}_2)}{r_{12}}\,d\mathbf{r}_1\,d\mathbf{r}_2$ are the \emph{exchange integrals}, a purely quantum mechanical effect.
\end{itemize}

\subsection{Exchange and the Fermi Hole}

The exchange integrals $K_{ij}$ arise from the antisymmetry requirement and have no classical analogue. Physically, same-spin electrons cannot occupy the same point in space (Pauli exclusion), creating a depletion of same-spin electron density around each electron known as the \emph{Fermi hole} or \emph{exchange hole}. Because same-spin electrons are kept further apart than the mean-field Coulomb picture would suggest, their repulsion is reduced. Exchange is always stabilizing: $K_{ij} \geq 0$, so the $-K_{ij}$ term lowers the energy.

Key properties of the exchange integrals:
\begin{itemize}
\item $K_{ij} > 0$ always; exchange \emph{lowers} the energy.
\item $K_{ij}$ is non-zero only for same-spin electron pairs.
\item Hartree--Fock theory treats exchange \emph{exactly} within the single-determinant ansatz.
\end{itemize}

\subsection{Electron Correlation}
\label{sec:correlation}

Correlation is, by definition, everything that the Hartree--Fock approximation misses:
\begin{equation}
E_\text{corr} = E_\text{exact} - E_\text{HF}.
\end{equation}
In the HF picture, each electron moves in the average (mean) field of all other electrons. It does not respond to the instantaneous positions of other electrons. In reality, electrons avoid each other dynamically, creating a \emph{Coulomb hole} around each electron in addition to the exchange hole.

Two distinct types of correlation are recognized:
\begin{description}
\item[Dynamic correlation] arises from the short-range avoidance of electron pairs. It is a smooth effect contributed by many electron configurations and is captured by methods such as M\o ller--Plesset perturbation theory (MP2) and coupled cluster theory (CCSD, CCSD(T)).
\item[Static (non-dynamic) correlation] arises when multiple electron configurations are nearly degenerate, and a single Slater determinant is qualitatively inadequate. Multi-reference methods (CASSCF, MRCI) are required for such cases.
\end{description}

The correlation energy is typically small in absolute terms---about 1\% of the total electronic energy---but enormous on a chemical scale. For water, $E_\text{corr} \approx \SI{-0.37}{E_h} \approx \SI{-230}{kcal/mol}$, while a typical chemical bond is worth \SIrange{50}{100}{kcal/mol}. Post-HF methods recover this energy at considerable computational cost: MP2 scales as $\mathcal{O}(N^5)$, CCSD as $\mathcal{O}(N^6)$, and CCSD(T)---the ``gold standard'' of quantum chemistry---as $\mathcal{O}(N^7)$, where $N$ is a measure of system size. For the molecules of interest to most chemists, these methods are frequently prohibitive. This motivates density functional theory.


% ============================================================================
\section{Density Functional Theory}
\label{sec:DFT}
% ============================================================================

\subsection{The Electron Density}

DFT adopts a fundamentally different strategy from wavefunction-based methods. Instead of the $N$-electron wavefunction $\Psi(\mathbf{r}_1, \mathbf{r}_2, \ldots, \mathbf{r}_N)$, which depends on $3N$ spatial coordinates, DFT works with the electron density $\rho(\mathbf{r})$, which depends on only three:
\begin{equation}
\rho(\mathbf{r}) = N\int |\Psi(\mathbf{r}, \mathbf{r}_2, \ldots, \mathbf{r}_N)|^2\,d\mathbf{r}_2\cdots d\mathbf{r}_N.
\label{eq:density}
\end{equation}
The density is non-negative, integrates to the total number of electrons ($\int \rho(\mathbf{r})\,d\mathbf{r} = N$), and---crucially---is a physical observable measurable by X-ray diffraction.

\subsection{The Hohenberg--Kohn Theorems}

The theoretical foundation of DFT rests on two theorems proved by Hohenberg and Kohn in 1964~\cite{HohenbergKohn1964}:

\begin{description}
\item[Theorem 1 (Existence):] The ground-state electron density $\rho_0(\mathbf{r})$ uniquely determines the external potential $v_\text{ext}(\mathbf{r})$ (up to a constant). Since $v_\text{ext}$ together with $N$ determines $\hat{H}$, all ground-state properties are functionals of $\rho_0$.

\item[Theorem 2 (Variational Principle):] For any trial density $\tilde{\rho}$ satisfying $\tilde{\rho}(\mathbf{r}) \geq 0$ and $\int\tilde{\rho}\,d\mathbf{r} = N$, the energy functional satisfies $E[\tilde{\rho}] \geq E_0$, with equality only for the true ground-state density.
\end{description}

The proof of Theorem 1 proceeds by \emph{reductio ad absurdum}. Assume two different external potentials $v_\text{ext}$ and $v_\text{ext}'$ (differing by more than a constant) give rise to the same ground-state density $\rho_0$. They define two different Hamiltonians $\hat{H}$ and $\hat{H}'$ with ground states $\Psi_0$ and $\Psi_0'$ and energies $E_0$ and $E_0'$. By the variational principle applied to $\hat{H}$:
\begin{equation}
E_0 < \braket{\Psi_0'|\hat{H}|\Psi_0'} = \braket{\Psi_0'|\hat{H}'|\Psi_0'} + \int [v_\text{ext}(\mathbf{r}) - v_\text{ext}'(\mathbf{r})]\rho_0(\mathbf{r})\,d\mathbf{r} = E_0' + \int [v_\text{ext} - v_\text{ext}']\rho_0\,d\mathbf{r}.
\end{equation}
By the same argument with primed and unprimed interchanged, $E_0' < E_0 + \int [v_\text{ext}' - v_\text{ext}]\rho_0\,d\mathbf{r}$. Adding the two inequalities gives $E_0 + E_0' < E_0' + E_0$, a contradiction. Therefore, the mapping $\rho_0 \to v_\text{ext}$ is unique.

These theorems are exact and profound: everything is encoded in $\rho(\mathbf{r})$. The difficulty is that the exact functional $E[\rho]$ is unknown.

\subsection{The Kohn--Sham Formulation}

The practical breakthrough came from Kohn and Sham~\cite{KohnSham1965}, who introduced a fictitious system of non-interacting electrons constrained to reproduce the exact ground-state density of the real interacting system. The Kohn--Sham (KS) energy functional is
\begin{equation}
E_\text{KS}[\rho] = T_s[\rho] + \int v_\text{ext}(\mathbf{r})\rho(\mathbf{r})\,d\mathbf{r} + \frac{1}{2}\iint\frac{\rho(\mathbf{r})\rho(\mathbf{r}')}{|\mathbf{r}-\mathbf{r}'|}\,d\mathbf{r}\,d\mathbf{r}' + E_\text{XC}[\rho],
\label{eq:EKS}
\end{equation}
where:
\begin{itemize}
\item $T_s[\rho]$ is the kinetic energy of the non-interacting KS electrons, computed exactly from the KS orbitals,
\item the second term is the electron--nuclear attraction,
\item the third term ($J[\rho]$) is the classical Coulomb repulsion,
\item $E_\text{XC}[\rho]$ is the exchange-correlation functional, which contains everything else: exchange energy, correlation energy, and the kinetic energy correction $T[\rho] - T_s[\rho]$.
\end{itemize}

Three of these four terms are computed exactly. The entire challenge of DFT reduces to approximating a single term, $E_\text{XC}[\rho]$.

\subsection{The Kohn--Sham Equations}

Minimizing $E_\text{KS}[\rho]$ subject to the constraint that the KS orbitals remain orthonormal yields one-electron equations:
\begin{equation}
\left[-\frac{1}{2}\nabla^2 + v_\text{eff}(\mathbf{r})\right]\phi_i(\mathbf{r}) = \varepsilon_i\,\phi_i(\mathbf{r}),
\label{eq:KS_equations}
\end{equation}
with the effective potential
\begin{equation}
v_\text{eff}(\mathbf{r}) = v_\text{ext}(\mathbf{r}) + \int\frac{\rho(\mathbf{r}')}{|\mathbf{r}-\mathbf{r}'|}\,d\mathbf{r}' + v_\text{XC}(\mathbf{r}),
\end{equation}
where $v_\text{XC}(\mathbf{r}) = \delta E_\text{XC}/\delta\rho(\mathbf{r})$ is the functional derivative of the exchange-correlation energy.

The density is reconstructed from the occupied KS orbitals:
\begin{equation}
\rho(\mathbf{r}) = \sum_{i=1}^{N}|\phi_i(\mathbf{r})|^2.
\end{equation}

Because $v_\text{eff}$ depends on $\rho$, which in turn depends on the orbitals $\{\phi_i\}$, the KS equations must be solved iteratively by the \emph{self-consistent field} (SCF) procedure:
\begin{enumerate}
\item Guess an initial density $\rho^{(0)}$.
\item Construct $v_\text{eff}$ from $\rho^{(0)}$.
\item Solve the KS equations to obtain orbitals $\{\phi_i^{(1)}\}$.
\item Compute $\rho^{(1)} = \sum_i |\phi_i^{(1)}|^2$.
\item If $|\rho^{(1)} - \rho^{(0)}|$ is below threshold, stop; otherwise, return to step 2.
\end{enumerate}

\subsection{Hartree--Fock vs.\ Kohn--Sham DFT}

Table~\ref{tab:HFvsDFT} summarizes the key differences.

\begin{table}[h]
\centering
\caption{Comparison of Hartree--Fock and Kohn--Sham DFT.}
\label{tab:HFvsDFT}
\begin{tabular}{lcc}
\toprule
& \textbf{Hartree--Fock} & \textbf{Kohn--Sham DFT} \\
\midrule
Central quantity & Wavefunction $\Psi$ & Density $\rho(\mathbf{r})$ \\
Exchange & Exact & Approximate ($E_\text{XC}$) \\
Correlation & Missing & Included in $E_\text{XC}$ \\
Cost & $\mathcal{O}(N^3$--$N^4)$ & $\mathcal{O}(N^3)$ \\
Systematic improvement & Yes (MP2, CCSD, \ldots) & No guarantee \\
\bottomrule
\end{tabular}
\end{table}

The principal advantage of DFT is that it includes correlation effects at nearly the same computational cost as HF. The principal disadvantage is that $E_\text{XC}$ is approximate, and there is no guaranteed route to systematic improvement.


% ============================================================================
\section{Jacob's Ladder of Density Functional Approximations}
\label{sec:jacobs_ladder}
% ============================================================================

Perdew introduced the metaphor of ``Jacob's Ladder''~\cite{Perdew2001} to organize the hierarchy of exchange-correlation functional approximations, with each successive rung incorporating additional information about the density and, in principle, increasing accuracy.

\subsection{Rung 1: Local Density Approximation (LDA)}

The LDA assumes that the exchange-correlation energy density at each point in space depends only on the local value of $\rho$:
\begin{equation}
E_\text{XC}^\text{LDA}[\rho] = \int \epsilon_\text{XC}^\text{hom}(\rho(\mathbf{r}))\,\rho(\mathbf{r})\,d\mathbf{r},
\end{equation}
where $\epsilon_\text{XC}^\text{hom}$ is the exchange-correlation energy per electron of the uniform electron gas at density $\rho$. The exchange part was derived analytically by Dirac, while the correlation part was parameterized by Vosko, Wilk, and Nusair from quantum Monte Carlo calculations. LDA works well for solids with slowly varying densities but overbinds molecules and gives poor barrier heights.

\subsection{Rung 2: Generalized Gradient Approximation (GGA)}

GGA functionals supplement $\rho$ with its gradient $\nabla\rho$:
\begin{equation}
E_\text{XC}^\text{GGA}[\rho] = \int f(\rho, \nabla\rho)\,d\mathbf{r}.
\end{equation}
Important examples include PBE~\cite{PBE1996} (derived from exact constraints, popular in solid-state chemistry) and BLYP (combining Becke's 1988 exchange with Lee--Yang--Parr correlation). GGA represents a substantial improvement over LDA for molecular properties.

\subsection{Rung 3: Meta-GGA}

Meta-GGA functionals further incorporate the kinetic energy density $\tau(\mathbf{r}) = \frac{1}{2}\sum_i |\nabla\phi_i(\mathbf{r})|^2$ and/or the Laplacian $\nabla^2\rho$. TPSS, SCAN~\cite{SCAN2015}, and r$^2$SCAN~\cite{r2SCAN2020} are prominent examples. SCAN, in particular, satisfies all 17 known exact constraints that a meta-GGA can satisfy.

\subsection{Rung 4: Hybrid Functionals}

Hybrid functionals mix a fraction of exact (Hartree--Fock) exchange into the DFT exchange:
\begin{equation}
E_\text{XC}^\text{hybrid} = a\,E_\text{X}^\text{HF} + (1-a)\,E_\text{X}^\text{DFA} + E_\text{C}^\text{DFA},
\end{equation}
where $a$ is the fraction of HF exchange. The most widely used functional in chemistry, B3LYP~\cite{Becke1993,Stephens1994}, uses $a = 0.20$ (20\% HF exchange). PBE0 uses $a = 0.25$. Range-separated hybrids such as $\omega$B97X-D vary $a$ with the interelectronic distance, using more HF exchange at long range; this improves the description of charge-transfer states and reaction barriers.

\subsection{Rung 5: Double Hybrids}

Double-hybrid functionals additionally incorporate a fraction of MP2-like correlation, e.g., B2PLYP~\cite{B2PLYP2006}. These are the most accurate DFT functionals for thermochemistry and reaction barriers but carry the $\mathcal{O}(N^5)$ cost of the MP2 component.

\subsection{Practical Recommendations}

Table~\ref{tab:functionals} provides guidance for common applications.

\begin{table}[h]
\centering
\caption{Recommended functionals for different applications.}
\label{tab:functionals}
\begin{tabular}{llcc}
\toprule
\textbf{Application} & \textbf{Recommended Functional} & \textbf{Rung} & \textbf{Typical Error} \\
\midrule
Organic thermochemistry & B3LYP-D3(BJ) or $\omega$B97X-D & 4 & \SIrange{1}{3}{kcal/mol} \\
Conformational energies & r$^2$SCAN-D4 or $\omega$B97X-D & 3--4 & \SIrange{0.5}{2}{kcal/mol} \\
Transition metals & PBE0-D3(BJ) or TPSSh & 4 & \SIrange{2}{5}{kcal/mol} \\
Non-covalent interactions & $\omega$B97X-D or B3LYP-D3(BJ) & 4 & \SIrange{0.2}{1}{kcal/mol} \\
High-accuracy benchmarks & DLPNO-CCSD(T) or B2PLYP-D3 & 5/WFT & ${<}\SI{1}{kcal/mol}$ \\
Solids and surfaces & PBE-D3 or r$^2$SCAN & 2--3 & varies \\
\bottomrule
\end{tabular}
\end{table}


% ============================================================================
\section{Dispersion Corrections}
\label{sec:dispersion}
% ============================================================================

Standard density functional approximations up to and including many hybrid functionals fail to capture London dispersion interactions---the attractive van der Waals forces arising from correlated instantaneous electron fluctuations. These forces are present between \emph{all} atoms and molecules, scale as $R^{-6}$ at long range, and can contribute several kcal/mol to relative energies. Neglecting dispersion leads to systematic errors in conformational energies, molecular crystal structures, binding energies of supramolecular complexes, and reaction barriers involving large substrates.

Several approaches to correcting this deficiency are in widespread use:

\begin{description}
\item[DFT-D3~\cite{GrimmeD3}:] Grimme's D3 correction adds pairwise $C_6/R^6$ (and optionally $C_8/R^8$) terms, with Becke--Johnson (BJ) damping~\cite{GrimmeD3BJ} to avoid double-counting at short range. The dispersion coefficients $C_6^{AB}$ depend on the coordination number, providing geometry-dependent polarizabilities. D3(BJ) is the most widely used correction and is recommended as a default.

\item[DFT-D4~\cite{GrimmeD4}:] An extension of D3 that makes the dispersion coefficients dependent on partial atomic charges, improving the description of polar and charged systems.

\item[XDM~\cite{BeckeJohnson2005,PriceXDM2023}:] The exchange-hole dipole moment model derives dispersion coefficients from the properties of the exchange hole rather than from external parameters. It provides an \emph{ab initio} route to $C_6$, $C_8$, and $C_{10}$ coefficients. In my own research, I have investigated the formal requirements for accurate dispersion modelling~\cite{PriceDispersion2021} and implemented the XDM model within the FHI-aims electronic structure code~\cite{PriceXDM2023}.

\item[Many-body dispersion (MBD)~\cite{TkatchenkoScheffler2012}:] Goes beyond pairwise additivity by treating the coupled fluctuating dipoles of all atoms simultaneously. Important for condensed phases and large molecular assemblies.
\end{description}

The practical message is straightforward: \emph{always} include a dispersion correction. In Gaussian, this requires only the keyword \texttt{EmpiricalDispersion=GD3BJ}; in ORCA, the keyword \texttt{D3BJ} or \texttt{D4}. The computational overhead is negligible.


% ============================================================================
\section{Basis Sets}
\label{sec:basis}
% ============================================================================

\subsection{Slater-Type vs.\ Gaussian-Type Orbitals}

The Kohn--Sham (or Hartree--Fock) orbitals must be expanded in a set of known functions:
\begin{equation}
\phi_i(\mathbf{r}) = \sum_\mu c_{\mu i}\,\chi_\mu(\mathbf{r}),
\end{equation}
where $\{\chi_\mu\}$ are the basis functions and $\{c_{\mu i}\}$ are the expansion coefficients determined by the SCF procedure.

Slater-type orbitals (STOs) have the correct exponential decay at long range and the correct cusp at the nucleus, but the required multi-center integrals are computationally intractable. Gaussian-type orbitals (GTOs, $\sim e^{-\alpha r^2}$) have the wrong functional form---they decay too quickly and lack the nuclear cusp---but all required integrals can be evaluated analytically. Modern basis sets use contracted GTOs: fixed linear combinations of ``primitive'' Gaussians that approximate the shape of STOs.

\subsection{Basis Set Families}

Three major families are in common use:

\begin{description}
\item[Pople basis sets] (e.g., 6-31G(d,p), 6-311+G(d,p)): The notation ``6-311+G(d,p)'' means the core is described by a contraction of 6 primitives; the valence is split into three sets (triple-$\zeta$) of 3, 1, and 1 primitives; ``+'' adds diffuse functions on heavy atoms; ``(d,p)'' adds $d$-type polarization on heavy atoms and $p$-type on hydrogen.

\item[Dunning correlation-consistent basis sets] (cc-pVDZ, cc-pVTZ, cc-pVQZ, \ldots): Designed to converge systematically toward the complete basis set (CBS) limit, making them ideal for extrapolation studies and wavefunction-based methods. The aug- prefix denotes diffuse augmentation.

\item[Karlsruhe def2 basis sets] (def2-SVP, def2-TZVP, def2-TZVPP, def2-QZVPP): Efficient, well-balanced sets that are my personal recommendation for routine DFT calculations. def2-TZVP provides an excellent cost-accuracy balance for most applications.
\end{description}

\subsection{Basis Set Superposition Error}

When computing binding energies of weakly bound complexes, an artifact called \emph{basis set superposition error} (BSSE) arises: each monomer ``borrows'' basis functions from its partner, artificially lowering the complex energy. The counterpoise correction of Boys and Bernardi~\cite{BoysBernardi1970} estimates and removes this error. BSSE diminishes with increasing basis set size and is largely negligible at the quadruple-$\zeta$ level.

The minimum standard for publishable DFT results is a triple-$\zeta$ basis with polarization functions (e.g., def2-TZVP, cc-pVTZ, or 6-311+G(d,p)).


% ============================================================================
\section{The Potential Energy Surface}
\label{sec:PES}
% ============================================================================

After the Born--Oppenheimer approximation, the total energy $E(\mathbf{R})$ is a function of the nuclear coordinates $\mathbf{R}$. For $M$ atoms, this function defines a surface in $3M-6$ dimensions ($3M-5$ for a linear molecule)---the potential energy surface (PES).

The key features of the PES are:
\begin{itemize}
\item \textbf{Minima}: stationary points where $\nabla E = \mathbf{0}$ and all eigenvalues of the Hessian are positive. These correspond to stable molecular structures---reactants, products, intermediates, and conformers.
\item \textbf{First-order saddle points}: stationary points with exactly one negative Hessian eigenvalue. These are transition states (discussed in the Lecture~2 notes).
\item \textbf{Gradient} $\mathbf{g} = \nabla E$: the force on the atoms (with opposite sign). Governs geometry optimization.
\item \textbf{Hessian} $\mathbf{H}$: the matrix of second derivatives, $H_{ij} = \partial^2 E/\partial R_i\partial R_j$. Governs vibrational analysis.
\end{itemize}

\subsection{Geometry Optimization}

A geometry optimization seeks the nearest minimum on the PES starting from an initial guess. The simplest approach, steepest descent, takes steps opposite to the gradient:
\begin{equation}
\mathbf{R}_{n+1} = \mathbf{R}_n - \alpha\,\mathbf{g}_n,
\end{equation}
where $\alpha$ is a step size. More efficient algorithms use second-derivative (Hessian) information. The quasi-Newton method updates an approximate Hessian at each step:
\begin{equation}
\mathbf{R}_{n+1} = \mathbf{R}_n - \mathbf{H}_n^{-1}\,\mathbf{g}_n.
\label{eq:newton_raphson}
\end{equation}
Convergence is declared when the maximum and RMS values of the gradient and the displacement fall below specified thresholds (typically $\SI{4.5e-4}{E_h/a_0}$ for the maximum force in Gaussian).


% ============================================================================
\section{The Harmonic Approximation}
\label{sec:harmonic}
% ============================================================================

DFT provides the electronic energy $E_\text{elec}(\mathbf{R})$, but thermodynamic quantities require knowledge of nuclear motion. The harmonic approximation provides the bridge.

\subsection{Taylor Expansion of the PES}

Near a minimum $\mathbf{R}_e$, the potential energy is expanded in a Taylor series:
\begin{equation}
V(\mathbf{R}) = V(\mathbf{R}_e) + \underbrace{\sum_i \frac{\partial V}{\partial R_i}\bigg|_{\mathbf{R}_e}}_{=\,0\text{ at minimum}} \Delta R_i + \frac{1}{2}\sum_{i,j}\frac{\partial^2 V}{\partial R_i\partial R_j}\bigg|_{\mathbf{R}_e}\Delta R_i\,\Delta R_j + \cdots
\label{eq:taylor}
\end{equation}
The first derivatives vanish at the minimum. The leading non-trivial contribution is quadratic, and the matrix of second derivatives is the Hessian:
\begin{equation}
H_{ij} = \frac{\partial^2 V}{\partial R_i\partial R_j}\bigg|_{\mathbf{R}_e}.
\end{equation}
For a diatomic, this reduces to the familiar harmonic potential $V(r) \approx \frac{1}{2}k(r-r_e)^2$, where $k = d^2V/dr^2|_{r_e}$ is the force constant.

\subsection{The Hessian Matrix and Normal Modes}

The Hessian is a $3M\times 3M$ symmetric matrix. To extract physically meaningful vibrational frequencies, we transform to mass-weighted coordinates:
\begin{equation}
q_i = \sqrt{m_i}\,\Delta R_i, \qquad \tilde{H}_{ij} = \frac{H_{ij}}{\sqrt{m_i\,m_j}}.
\end{equation}
Diagonalizing the mass-weighted Hessian,
\begin{equation}
\tilde{\mathbf{H}}\,\mathbf{L} = \mathbf{L}\,\boldsymbol{\Lambda},
\end{equation}
yields eigenvalues $\{\lambda_k\}$ and eigenvectors $\{\mathbf{l}_k\}$. The eigenvalues give the vibrational frequencies:
\begin{equation}
\nu_k = \frac{1}{2\pi}\sqrt{\lambda_k}.
\label{eq:freq}
\end{equation}
The eigenvectors are the \emph{normal modes}---independent patterns of atomic vibration in which all atoms oscillate at the same frequency and pass through equilibrium simultaneously.

Of the $3M$ eigenvalues, six (five for a linear molecule) are zero or near-zero, corresponding to the three translations and three (two) rotations that do not change the potential energy. The remaining $3M-6$ ($3M-5$) eigenvalues are the vibrational modes.

At a minimum, all vibrational eigenvalues are positive (real frequencies). A negative eigenvalue indicates an imaginary frequency and signals that the structure is not a minimum---it may be a transition state (one imaginary frequency) or a higher-order saddle point.

\subsection{Quantum Harmonic Oscillator}

Each normal mode $k$ behaves as an independent quantum harmonic oscillator with energy levels
\begin{equation}
\varepsilon_{k,n} = \left(n + \tfrac{1}{2}\right)h\nu_k, \qquad n = 0, 1, 2, \ldots
\label{eq:QHO}
\end{equation}
The most important feature is the \emph{zero-point energy} (ZPE): even at $T = 0$, the energy is $\frac{1}{2}h\nu_k$, not zero. This is a purely quantum mechanical effect---the Heisenberg uncertainty principle forbids the oscillator from resting at the potential minimum. The total ZPE is
\begin{equation}
E_\text{ZPE} = \frac{1}{2}\sum_{k=1}^{3M-6}h\nu_k.
\label{eq:ZPE}
\end{equation}
For a medium-sized organic molecule, $E_\text{ZPE}$ is typically tens of kcal/mol and \emph{must} be included in any comparison of molecular energies.


% ============================================================================
\section{Statistical Mechanics Connection}
\label{sec:statmech}
% ============================================================================

The machinery of statistical mechanics---developed in the earlier weeks of this course---converts microscopic energy levels into macroscopic thermodynamic quantities. The canonical partition function is
\begin{equation}
Q(N,V,T) = \sum_j e^{-E_j/k_BT},
\label{eq:partition}
\end{equation}
and for $N$ indistinguishable, independent molecules, $Q = q^N/N!$, where $q$ is the molecular partition function.

\subsection{Factorization of the Molecular Partition Function}

Under the Born--Oppenheimer, rigid-rotor, and harmonic-oscillator (RRHO) approximations, the molecular partition function factorizes:
\begin{equation}
\boxed{q = q_\text{trans}\cdot q_\text{rot}\cdot q_\text{vib}\cdot q_\text{elec}.}
\label{eq:q_factored}
\end{equation}
Each factor requires different information from the DFT calculation.

\subsection{Translational Partition Function}
\label{sec:qtrans}

For an ideal gas of mass $m$ in a volume $V$:
\begin{equation}
q_\text{trans} = \left(\frac{2\pi m k_BT}{h^2}\right)^{3/2}V = \frac{V}{\Lambda^3},
\label{eq:qtrans}
\end{equation}
where
\begin{equation}
\Lambda = \frac{h}{\sqrt{2\pi m k_BT}}
\label{eq:deBroglie}
\end{equation}
is the thermal de Broglie wavelength. For standard-state thermodynamics at pressure $P^\circ$, replace $V$ by $k_BT/P^\circ$.

At room temperature, $q_\text{trans} \sim 10^{30}$, making it by far the largest contribution. The translational entropy is given by the Sackur--Tetrode equation:
\begin{equation}
S_\text{trans} = R\left[\ln\!\left(\frac{q_\text{trans}}{N}\right) + \frac{5}{2}\right].
\label{eq:Sackur_Tetrode}
\end{equation}

\subsection{Rotational Partition Function}

For a nonlinear polyatomic molecule in the high-temperature (classical) limit:
\begin{equation}
q_\text{rot} = \frac{\sqrt{\pi}}{\sigma}\left(\frac{8\pi^2 k_BT}{h^2}\right)^{3/2}\!\sqrt{I_AI_BI_C},
\label{eq:qrot}
\end{equation}
where $I_A$, $I_B$, $I_C$ are the principal moments of inertia (from the DFT-optimized geometry) and $\sigma$ is the rotational symmetry number (e.g., $\sigma = 2$ for \ce{H2O}, $\sigma = 12$ for \ce{CH4} and \ce{C6H6}). For a linear molecule,
\begin{equation}
q_\text{rot}^\text{linear} = \frac{8\pi^2 I k_BT}{\sigma h^2}.
\end{equation}

The rotational entropy for a nonlinear molecule is
\begin{equation}
S_\text{rot} = R\left[\ln q_\text{rot} + \frac{3}{2}\right].
\end{equation}

\subsection{Vibrational Partition Function}

For a harmonic oscillator with frequency $\nu_k$, the partition function has the closed form
\begin{equation}
q_{\text{vib},k} = \frac{1}{1 - e^{-\Theta_{V,k}/T}},
\label{eq:qvib}
\end{equation}
where $\Theta_{V,k} = h\nu_k/k_B$ is the \emph{vibrational temperature}. The total vibrational partition function is a product over all modes:
\begin{equation}
q_\text{vib} = \prod_{k=1}^{3M-6}\frac{1}{1-e^{-\Theta_{V,k}/T}}.
\end{equation}
This expression uses the ``zero-point exclusive'' convention, measuring energies from the ZPE level. In the ``bottom-of-the-well'' convention used by most quantum chemistry programs:
\begin{equation}
q_\text{vib}^\text{BOW} = \prod_{k=1}^{3M-6}\frac{e^{-\Theta_{V,k}/2T}}{1-e^{-\Theta_{V,k}/T}}.
\end{equation}

Two limiting cases are instructive:
\begin{itemize}
\item $T \ll \Theta_V$: the mode is ``frozen out''; only the ground state is populated, $q_{\text{vib},k} \approx 1$, and the mode contributes negligibly to the entropy.
\item $T \gg \Theta_V$: the classical limit, $q_{\text{vib},k} \approx T/\Theta_V$, and the mode is fully thermally active.
\end{itemize}

For reference: a C--H stretch at $\sim\!\SI{3000}{cm^{-1}}$ has $\Theta_V \approx \SI{4300}{K}$ (frozen at room temperature), while a low-frequency torsion at $\sim\!\SI{100}{cm^{-1}}$ has $\Theta_V \approx \SI{144}{K}$ (fully classical at room temperature).

The vibrational contributions to the thermal energy and entropy are:
\begin{align}
U_\text{vib} &= R\sum_k \frac{\Theta_{V,k}}{e^{\Theta_{V,k}/T}-1}, \label{eq:Uvib}\\
S_\text{vib} &= R\sum_k\left[\frac{\Theta_{V,k}/T}{e^{\Theta_{V,k}/T}-1} - \ln\!\left(1-e^{-\Theta_{V,k}/T}\right)\right]. \label{eq:Svib}
\end{align}

\subsection{Electronic Partition Function}

For most closed-shell molecules at room temperature, electronic excitation energies are much larger than $k_BT \approx \SI{0.026}{eV}$ at \SI{298}{K}, and only the ground state is populated:
\begin{equation}
q_\text{elec} \approx g_0,
\end{equation}
where $g_0$ is the ground-state degeneracy (1 for a singlet, $2S+1$ for open-shell species). The electronic contribution to the entropy is $S_\text{elec} = R\ln g_0$.


% ============================================================================
\section{From DFT to Gibbs Free Energy}
\label{sec:G_from_DFT}
% ============================================================================

Collecting all contributions, the Gibbs free energy from a DFT frequency calculation is:
\begin{equation}
\boxed{G = E_\text{elec} + E_\text{ZPE} + H_\text{thermal} - TS,}
\label{eq:G}
\end{equation}
where
\begin{align}
E_\text{ZPE} &= \frac{1}{2}\sum_k h\nu_k, \\
H_\text{thermal} &= U_\text{trans} + U_\text{rot} + U_\text{vib} + k_BT \nonumber\\
&= \frac{3}{2}k_BT + \frac{3}{2}k_BT + \sum_k\frac{h\nu_k}{e^{h\nu_k/k_BT}-1} + k_BT, \label{eq:Hthermal}\\
S &= S_\text{trans} + S_\text{rot} + S_\text{vib} + S_\text{elec}. \label{eq:Stotal}
\end{align}
The $k_BT$ term in $H_\text{thermal}$ is the $PV = nRT$ contribution converting internal energy to enthalpy at constant pressure.

The electronic energy $E_\text{elec}$ dominates in absolute magnitude---typically thousands of hartrees for a medium-sized molecule. But the thermal corrections (ZPE, enthalpy, and entropy terms), though small in absolute value (millihartrees), are what determine the \emph{relative} energies that govern chemistry. When comparing two species, the large electronic energies largely cancel, and the thermal corrections tip the balance.


% ============================================================================
\section{Practical Considerations}
\label{sec:practical}
% ============================================================================

\subsection{Reading Gaussian and ORCA Output}

After a frequency calculation, Gaussian prints a thermochemistry summary of the form:
\begin{verbatim}
  Zero-point correction=                   0.044749 (Hartree/Particle)
  Thermal correction to Energy=            0.047545
  Thermal correction to Enthalpy=          0.048489
  Thermal correction to Gibbs Free Energy= 0.021178
  Sum of electronic and zero-point Energies=     -76.377939
  Sum of electronic and thermal Energies=        -76.375143
  Sum of electronic and thermal Enthalpies=      -76.374199
  Sum of electronic and thermal Free Energies=   -76.401510
\end{verbatim}

For computing reaction free energies, the key line is ``Sum of electronic and thermal Free Energies''---this is $G$ for the species. For a reaction \ce{A -> B}:
\begin{equation}
\Delta G^\circ = G(B) - G(A).
\end{equation}
The equilibrium constant follows from
\begin{equation}
K = e^{-\Delta G^\circ/RT}.
\end{equation}
By default, Gaussian evaluates thermochemistry at $T = \SI{298.15}{K}$ and $P = \SI{1}{atm}$; these can be changed with the \texttt{Temperature} and \texttt{Pressure} keywords.

In ORCA, the equivalent information appears in the ``THERMOCHEMISTRY'' block near the end of the output, labeled ``Final Gibbs free enthalpy.''

\subsection{Hierarchy of Energy Corrections}

It is essential to distinguish the hierarchy of energy corrections:
\begin{align}
\Delta E &= E_\text{elec}(\text{products}) - E_\text{elec}(\text{reactants}), \label{eq:DeltaE}\\
\Delta E_0 &= \Delta E + \Delta E_\text{ZPE}, \label{eq:DeltaE0}\\
\Delta H &= \Delta E_0 + \Delta H_\text{thermal}, \label{eq:DeltaH}\\
\Delta G &= \Delta H - T\Delta S. \label{eq:DeltaG}
\end{align}
These are \emph{all different numbers}. For meaningful comparison with experiment, $\Delta G$ is required for equilibria and $\Delta G^\ddagger$ for kinetics (Lecture~2). The zero-point energy correction alone can shift relative energies by several kcal/mol.

\subsection{Quasi-Harmonic Corrections}
\label{sec:quasi_harmonic}

The harmonic approximation breaks down for low-frequency modes (below $\sim\!\SI{100}{cm^{-1}}$), which correspond to large-amplitude motions such as internal rotations and ring puckering. These modes contribute disproportionately to the vibrational entropy through Eq.~\eqref{eq:Svib} because $S_{\text{vib},k}$ diverges as $\nu_k \to 0$, a clearly unphysical result.

Two widely used corrections address this:

\begin{description}
\item[Truhlar's quasi-harmonic approximation~\cite{TruhlarQH}:] All vibrational frequencies below \SI{100}{cm^{-1}} are raised to \SI{100}{cm^{-1}} for the purpose of computing the entropy. The ZPE continues to use the original frequencies.

\item[Grimme's quasi-RRHO~\cite{GrimmeQRRHO2012}:] Below a cutoff frequency ($\omega_0 \approx \SI{100}{cm^{-1}}$), the vibrational contribution to the entropy is smoothly interpolated between the harmonic oscillator model and a free-rotor model. Specifically, for mode $k$ with frequency $\nu_k$:
\begin{equation}
S_k = w(\nu_k)\,S_k^\text{HO} + [1 - w(\nu_k)]\,S_k^\text{FR},
\end{equation}
where $w(\nu_k)$ is a switching function that equals 1 for $\nu_k \gg \omega_0$ and 0 for $\nu_k \ll \omega_0$, and $S_k^\text{FR}$ is the free-rotor entropy. This is more physically motivated than the hard cutoff and is the default in ORCA thermochemistry.
\end{description}

\subsection{ZPE Scaling Factors}

Computed harmonic frequencies systematically overestimate experimental fundamental frequencies because they neglect anharmonicity. Empirical scaling factors (e.g., 0.9804 for B3LYP/6-31G(d), 0.9854 for B3LYP/def2-TZVP) are applied to bring computed ZPEs into better agreement with experiment. These factors are tabulated for many functional/basis combinations in the NIST CCCBDB~\cite{CCCBDB}.


% ============================================================================
\section{Worked Example: Reaction Free Energy of the Water--Gas Shift Reaction}
\label{sec:worked}
% ============================================================================

Consider the water--gas shift reaction:
\begin{equation}
\ce{CO(g) + H2O(g) -> CO2(g) + H2(g)}.
\end{equation}

A complete DFT study at the B3LYP-D3(BJ)/def2-TZVP level proceeds as follows:

\begin{enumerate}
\item \textbf{Geometry optimization + frequency calculation} for each species: \ce{CO}, \ce{H2O}, \ce{CO2}, \ce{H2}. Confirm all have zero imaginary frequencies (true minima).

\item \textbf{Read the Gibbs free energies} $G$ from each output file (``Sum of electronic and thermal Free Energies'' at \SI{298.15}{K}, \SI{1}{atm}).

\item \textbf{Compute} $\Delta G^\circ_\text{rxn}$:
\begin{equation}
\Delta G^\circ_\text{rxn} = \bigl[G(\ce{CO2}) + G(\ce{H2})\bigr] - \bigl[G(\ce{CO}) + G(\ce{H2O})\bigr].
\end{equation}

\item \textbf{Compute the equilibrium constant}:
\begin{equation}
K = e^{-\Delta G^\circ_\text{rxn}/RT}.
\end{equation}
\end{enumerate}

Suppose the DFT calculation yields $\Delta G^\circ_\text{rxn} = \SI{-6.8}{kcal/mol}$ at \SI{298}{K}. Then:
\begin{equation}
K = e^{-(-6800)/(1.987 \times 298)} = e^{11.5} \approx 1.0 \times 10^5.
\end{equation}
This indicates strong thermodynamic favorability, consistent with the known exergonicity of the water--gas shift reaction.

To illustrate the importance of thermal corrections, the individual contributions might break down as:
\begin{center}
\begin{tabular}{lS[table-format=-2.1]}
\toprule
\textbf{Contribution} & {$\Delta$ (kcal/mol)} \\
\midrule
$\Delta E_\text{elec}$ & -5.2 \\
$\Delta E_\text{ZPE}$ & -1.8 \\
$\Delta H_\text{thermal}$ & +0.4 \\
$-T\Delta S$ & -0.2 \\
\midrule
$\Delta G^\circ$ & -6.8 \\
\bottomrule
\end{tabular}
\end{center}

The ZPE correction of $\SI{-1.8}{kcal/mol}$ is substantial---omitting it would underestimate the driving force by more than 25\%. This underscores the necessity of performing frequency calculations, not just geometry optimizations.


% ============================================================================
\section{Summary}
\label{sec:summary}
% ============================================================================

The path from the molecular Schr\"{o}dinger equation to computed thermodynamic properties proceeds through a series of well-defined steps:

\begin{enumerate}
\item The \textbf{Born--Oppenheimer approximation} separates the electronic and nuclear problems.
\item \textbf{Kohn--Sham DFT} solves the electronic problem at each nuclear geometry, yielding $E_\text{elec}(\mathbf{R})$.
\item \textbf{Geometry optimization} locates minima on the PES.
\item The \textbf{harmonic approximation} (Taylor expansion of the PES at the minimum) provides vibrational frequencies from the Hessian.
\item \textbf{Statistical mechanics} converts these frequencies (together with the molecular mass and geometry) into partition functions: $q = q_\text{trans}\cdot q_\text{rot}\cdot q_\text{vib}\cdot q_\text{elec}$.
\item The thermodynamic properties follow: $G = E_\text{elec} + E_\text{ZPE} + H_\text{thermal} - TS$.
\end{enumerate}

A single DFT frequency calculation provides every quantity needed to evaluate all four partition function factors and, from them, $\Delta G$, $\Delta H$, $\Delta S$, $K_\text{eq}$, and---as shown in the Lecture~2 notes---rate constants and selectivities.


% ============================================================================
\section*{References}
% ============================================================================
\addcontentsline{toc}{section}{References}

\begin{enumerate}[label={[\arabic*]},leftmargin=*]
\bibitem{SzaboOstlund} A.\ Szabo and N.\ S.\ Ostlund, \textit{Modern Quantum Chemistry: Introduction to Advanced Electronic Structure Theory} (Dover, 1996).

\bibitem{Jensen} F.\ Jensen, \textit{Introduction to Computational Chemistry}, 3rd ed.\ (Wiley, 2017).

\bibitem{Cramer} C.\ J.\ Cramer, \textit{Essentials of Computational Chemistry: Theories and Models}, 2nd ed.\ (Wiley, 2004).

\bibitem{KochHolthausen} W.\ Koch and M.\ C.\ Holthausen, \textit{A Chemist's Guide to Density Functional Theory}, 2nd ed.\ (Wiley-VCH, 2001).

\bibitem{ParrYang} R.\ G.\ Parr and W.\ Yang, \textit{Density-Functional Theory of Atoms and Molecules} (Oxford University Press, 1989).

\bibitem{BornOppenheimer1927} M.\ Born and R.\ Oppenheimer, Ann.\ Phys.\ \textbf{389}, 457 (1927).

\bibitem{HohenbergKohn1964} P.\ Hohenberg and W.\ Kohn, Phys.\ Rev.\ \textbf{136}, B864 (1964).

\bibitem{KohnSham1965} W.\ Kohn and L.\ J.\ Sham, Phys.\ Rev.\ \textbf{140}, A1133 (1965).

\bibitem{Perdew2001} J.\ P.\ Perdew and K.\ Schmidt, in \textit{Density Functional Theory and Its Application to Materials}, AIP Conf.\ Proc.\ \textbf{577}, 1 (2001).

\bibitem{PBE1996} J.\ P.\ Perdew, K.\ Burke, and M.\ Ernzerhof, Phys.\ Rev.\ Lett.\ \textbf{77}, 3865 (1996).

\bibitem{Becke1993} A.\ D.\ Becke, J.\ Chem.\ Phys.\ \textbf{98}, 5648 (1993).

\bibitem{Stephens1994} P.\ J.\ Stephens, F.\ J.\ Devlin, C.\ F.\ Chabalowski, and M.\ J.\ Frisch, J.\ Phys.\ Chem.\ \textbf{98}, 11623 (1994).

\bibitem{SCAN2015} J.\ Sun, A.\ Ruzsinszky, and J.\ P.\ Perdew, Phys.\ Rev.\ Lett.\ \textbf{115}, 036402 (2015).

\bibitem{r2SCAN2020} J.\ W.\ Furness, A.\ D.\ Kaplan, J.\ Ning, J.\ P.\ Perdew, and J.\ Sun, J.\ Phys.\ Chem.\ Lett.\ \textbf{11}, 8208 (2020).

\bibitem{B2PLYP2006} S.\ Grimme, J.\ Chem.\ Phys.\ \textbf{124}, 034108 (2006).

\bibitem{GrimmeD3} S.\ Grimme, J.\ Antony, S.\ Ehrlich, and H.\ Krieg, J.\ Chem.\ Phys.\ \textbf{132}, 154104 (2010).

\bibitem{GrimmeD3BJ} S.\ Grimme, S.\ Ehrlich, and L.\ Goerigk, J.\ Comput.\ Chem.\ \textbf{32}, 1456 (2011).

\bibitem{GrimmeD4} E.\ Caldeweyher, C.\ Bannwarth, and S.\ Grimme, J.\ Chem.\ Phys.\ \textbf{147}, 034112 (2017).

\bibitem{BeckeJohnson2005} A.\ D.\ Becke and E.\ R.\ Johnson, J.\ Chem.\ Phys.\ \textbf{122}, 154104 (2005).

\bibitem{PriceDispersion2021} A.\ J.\ Price, K.\ R.\ Bryenton, and E.\ R.\ Johnson, J.\ Chem.\ Phys.\ \textbf{154}, 230902 (2021).

\bibitem{PriceXDM2023} A.\ J.\ Price, A.\ Otero-de-la-Roza, and E.\ R.\ Johnson, Chem.\ Sci.\ \textbf{14}, 1252 (2023).

\bibitem{TkatchenkoScheffler2012} A.\ Tkatchenko, R.\ A.\ DiStasio Jr., R.\ Car, and M.\ Scheffler, Phys.\ Rev.\ Lett.\ \textbf{108}, 236402 (2012).

\bibitem{BoysBernardi1970} S.\ F.\ Boys and F.\ Bernardi, Mol.\ Phys.\ \textbf{19}, 553 (1970).

\bibitem{TruhlarQH} R.\ F.\ Ribeiro, A.\ V.\ Marenich, C.\ J.\ Cramer, and D.\ G.\ Truhlar, J.\ Phys.\ Chem.\ B \textbf{115}, 14556 (2011).

\bibitem{GrimmeQRRHO2012} S.\ Grimme, Chem.\ Eur.\ J.\ \textbf{18}, 9955 (2012).

\bibitem{CCCBDB} NIST Computational Chemistry Comparison and Benchmark Database, \url{https://cccbdb.nist.gov/}.
\end{enumerate}

\end{document}
